\documentclass[aspectratio=169,9pt]{beamer}
\mode<presentation> 
\usepackage[utf8]{inputenc}
\usepackage{sansmathaccent} %Für equations
\usepackage{amsmath}
\usepackage{tikz}
\usepackage{multicol}
\usepackage{cancel}
\usepackage{relsize}
\usepackage{graphicx}
\usepackage[table]{xcolor}
\usepackage{fontawesome5}
\usepackage{enumitem}
\usepackage{setspace}
\usepackage[labelfont={bf},font={color=black!40,footnotesize},tablename={Tabelle},figurename={Abbildung}]{caption}
\usetikzlibrary{positioning, arrows.meta, calc, matrix}
\usefonttheme[onlymath]{serif}
\usepackage{booktabs}
%\makeatletter
%\makeatother

\usetikzlibrary{patterns, arrows.meta, calc}
\usepackage{pgfplots}
\pgfplotsset{compat=1.18}
\usetikzlibrary{intersections, pgfplots.fillbetween}\usetikzlibrary{intersections, pgfplots.fillbetween}
\rowcolors{2}{miugray!30}{miugray!10}

\usetheme{Madrid}  %% Themenwahl
\useoutertheme{split} % Alternatively: miniframes, infolines, split
%\useinnertheme{circles}


\setbeamertemplate{itemize item}{\color{miugray}$\blacksquare$}
\setbeamertemplate{itemize subitem}{\color{uniblau}$\bullet$}



%\newcommand{\metermilli}{\mskip3mu \text{mm}}
\newcommand{\cm}{\mskip3mu \text{cm}}
\newcommand{\m}{ \text{m}}
\newcommand{\lite}{ \text{l}}
\newcommand{\kg}{ \text{kg}}
\newcommand{\g}{ \text{g}}
\newcommand{\km}{ \text{km}}
\newcommand{\h}{ \text{h}}
\newcommand{\s}{ \text{s}}
\newcommand{\tonne}{ \text{t}}


\newcommand{\RA}{$\rightarrow$~}
\newcommand{\pp}{\par \vspace*{2mm}}
\newcommand{\ppbig}{\par \vspace*{4mm}}

\definecolor{miudiutex}{HTML}{1d2926}
\definecolor{miugray}{HTML}{b4cccf}
\definecolor{miugray2}{HTML}{4d7365}
\definecolor{white2}{HTML}{f5f7f8}
\definecolor{red}{HTML}{cf1f5c}
\definecolor{green}{HTML}{00A36C}
\definecolor{delphinidin}{HTML}{315794}
\definecolor{uniblau}{HTML}{004e9f}
\definecolor{unigelb}{HTML}{fcba00}
\definecolor{unigrau}{HTML}{909085}
\definecolor{pelargonidin}{HTML}{b33807}
\definecolor{cyanidin}{HTML}{ba0746}
\definecolor{paeonidin}{HTML}{d15ca6}
\definecolor{petunidin}{HTML}{b56fed}
\definecolor{malvidin}{HTML}{8a2d8c}

%\setbeamercolor{palette primary}{bg=pelargonidin!35,fg=white}
%\setbeamercolor{palette secondary}{bg=delphinidin!45,fg=white}
\setbeamercolor{palette primary}{bg=uniblau!70,fg=white}
\setbeamercolor{palette secondary}{bg=uniblau!45,fg=white}
\setbeamercolor{palette tertiary}{bg=unigelb!70,fg=unigrau}
\setbeamercolor{palette quaternary}{bg=miugray,fg=white}
\setbeamercolor{structure}{fg=black} % itemize, enumerate, etc
\setbeamercolor{section in toc}{fg=red} % TOC sections

\setbeamersize{text margin left=1cm,text margin right=1cm}


\useinnertheme{tcolorbox} 
\newcommand{\goodtk}[2]{
\begin{tcolorbox}[
	enhanced,
	colframe=uniblau!75,
	sharp corners,
	boxsep=5pt,
	title=\bfseries \sffamily #1,
	colback=white,
	coltitle=black,
	attach boxed title to top text left={yshift=-\tcboxedtitleheight/2},
	boxed title style={colback=white,colframe=white}
]
	#2
	\end{tcolorbox}
}
	
\newcommand{\antwort}[1]{
\begin{tcolorbox}[
	enhanced,
	colframe=unigrau!65,
	sharp corners,
	boxsep=1pt,
	colback=white,
	coltitle=black,
]
	#1
	\end{tcolorbox}
}


\title{12. Übung}
\author{PD Dr. Helene Loos}
\date{\today}

\begin{document}
%\maketitle
\small



\begin{frame}
\begin{center}

\LARGE Übung 12  \par \normalsize
Prozessbezogene Lebensmitteltechnologie
\end{center}
\small
\end{frame}



\begin{frame}


\begin{minipage}{.5\textwidth}


{\color{miugray}$\blacksquare ~$} \textbf{Aufgabe 1} \par \vspace*{2mm}
Aus einem Orangensaftkonzentrat mit 30 °Brix soll ein Orangensaft mit 11,5 °Brix hergestellt
werden. Wieviel kg Saft kann hergestellt werden, wenn Sie 5 kg Orangensaftkonzentrat zur
Verfügung haben?

\end{minipage}%
\hspace*{4mm}
\begin{minipage}{.5\textwidth}
\antwort{
\uncover<2->{$ m_1 = 5~\text{kg};~m_2 = m_1 + m_3 ;~m_3 = ~?~;	\\ $ } 
\uncover<3->{$~x_1=30~\text{°};~x_2=11,5~\text{°};~x_3=0~\text{°}$}

\begin{align*}
\uncover<4->{m_1 \cdot x_1 + m_3 \cdot x_3 &= m_2 \cdot x_2 \\}
\uncover<5->{m_1 \cdot x_1 + m_3 \cdot x_3 &= m_1 \cdot x_2 + m_3 \cdot x_2 \\}
\uncover<6->{m_1 \cdot x_1 + m_3 \cdot x_3 - m_1 \cdot x_2 &= m_3 \cdot x_2 \\}
\uncover<7->{m_1 \cdot x_1 - m_1 \cdot x_2 &= m_3 \cdot x_2 - m_3 \cdot x_3\\}
\uncover<8->{m_1 \cdot x_1 - m_1 \cdot x_2 &= m_3 \cdot (x_2 - x_3)\\}
\uncover<9->{\frac{m_1 \cdot x_1 - m_1 \cdot x_2}{x_2 - x_3} &= m_3 = 8,04}
\end{align*}

\uncover<10->{Wir können also zu unseren vorhandenen 5 kg Orangensaftkonzentrat weitere 8,04 kg Wasser hinzugeben um einen °Brix von 11,5 zu erreichen. Dementsprechend erhalten wir 13,04 kg Orangensaft.}

}
\end{minipage}%


\end{frame}


\begin{frame}
\begin{minipage}{.4\textwidth}


{\color{miugray}$\blacksquare ~$} \textbf{Aufgabe 2} \par \vspace*{2mm}
\footnotesize
Ein Lebensmittel ist in Stickstoff-Schutzgasatmosphäre verpackt. Aufgrund eines Fehlers in
der Siegelnaht kommt es zur Diffusion von Sauerstoff und Wasserdampf aus der
Umgebungsluft in die Verpackung durch eine Pore (Porenradius: 12,5 µm, Porenlänge: 40
µm). Berechnen Sie den Stoffstrom für Sauerstoff und den Stoffstrom für Wasserdampf unter
der Annahme, dass die Konzentrationsdifferenzen von Sauerstoff und Wasserdampf 0,191 cm$^3$ (STP)/cm$^3$ und 12,91 µg/cm$^3$ betragen. Nehmen Sie für beide Gase einen Diffusionskoeffizienten von 0,243 cm$^2$/s an. 
\end{minipage}%
\hspace*{4mm}
\begin{minipage}{.6\textwidth}
\antwort{
\footnotesize
\textbf{Sauerstoff:}
\begin{align*}
\uncover<2->{\dot{n} &= D \cdot A \cdot \frac{\Delta c}{l} \\}
		\uncover<3->{&= 0,243 ~\text{cm}^2/\text{s} \cdot (\pi \cdot (1,25 \cdot 10^{-3}~\text{cm})^2) \cdot \frac{0,191 ~\text{cm}^3 \text{/cm$^3$}}{4 \cdot 10^{-3}~\text{cm} } \\}
		\uncover<4->{&= 5,69 \cdot 10^{-5} ~\text{cm$^3$/s}}
\end{align*}



\textbf{Wasser:}
\begin{align*}
\uncover<5->{\dot{n} &= D \cdot A \cdot \frac{\Delta c}{l} \\}
		\uncover<6->{&= 0,243 ~\text{cm}^2/\text{s} \cdot (\pi \cdot (1,25 \cdot 10^{-3}~\text{cm})^2) \cdot \frac{1,291 \cdot 10^{-5} ~\text{g}\text{/cm$^3$}}{4 \cdot 10^{-3}~\text{cm} } \\}
		\uncover<7->{&= 3,85 \cdot 10^{-9} ~\text{g/s}}
\end{align*}
}
\end{minipage}%


\end{frame}



\begin{frame}
\begin{minipage}{.4\textwidth}
%\normalsize

{\color{miugray}$\blacksquare ~$} \textbf{Aufgabe 3} \par \vspace*{2mm}
Ein Lebensmittel ist mit einem sauerstoff- und wasserdampfundurchlässigen Packstoff
verpackt, in dem jedoch eine Pore (Radius: 12,5 µm) ist. Die Pore ist durch eine
Kunststoffschicht mit der Dicke 40 µm abgedeckt. Berechnen Sie den Stoffstrom für
Sauerstoff und Wasserdampf, wenn der Partialdruckunterschied p$_{O2}$ 0,21 bar und der
Partialdruckunterschied p$_{H2O}$ 14,07 mbar beträgt. Der Permeationskoeffizient für Sauerstoff ist mit 628 (cm$^3$ (STP) $\cdot$ 100 µm)/(m$^2$ $\cdot$ d $\cdot$ bar) gegeben. Der Permeationskoeffizient für Wasserdampf beträgt (16,72 g $\cdot$ 100 µm)/(m$^2$ $\cdot$ d $\cdot$ bar).
\end{minipage}%
\hspace*{4mm}
\begin{minipage}{.6\textwidth}
\antwort{
\footnotesize
\uncover<2->{\begin{align*}
J &= P \cdot A \cdot \frac{\Delta p}{l}
\end{align*}}

\pp 
\textbf{Sauerstoff:}
\begin{align*}
\uncover<3->{J &= P \cdot A \cdot \frac{\Delta p}{l} \\}
 \uncover<4->{&= \frac{628 \cdot 10^{-6}~\text{m}^3 \cdot 10^{-4} ~\text{m}}{~\text{m}^2 \cdot ~\text{d} \cdot ~\text{bar}} \cdot \pi  \left( 1,25 \cdot 10^{-5}~\text{m} \right)^2  \cdot \frac{0,21 ~\text{bar}}{4 \cdot 10^{-5} ~\text{m}} \\}
  \uncover<5->{&= 1,62 \cdot 10^{-13} ~\text{m}^3 / ~\text{d}}
\end{align*}
\pp 
\textbf{Wasser:}
\begin{align*}
\uncover<6->{J &= P \cdot A \cdot \frac{\Delta p}{l} \\}
 \uncover<7->{&= \frac{16,72 ~\text{g}  \cdot 10^{-4} ~\text{m}}{~\text{m}^2 \cdot ~\text{d} \cdot ~\text{bar}} \cdot \pi  \left( 1,25 \cdot 10^{-5}~\text{m} \right)^2  \cdot \frac{14,07 \cdot 10^{-3} ~\text{bar}}{4 \cdot 10^{-5} ~\text{m}} \\}
  \uncover<8->{&= 2,88 \cdot 10^{-10} ~\text{g} / ~\text{d}}
\end{align*}

}
\end{minipage}%
\end{frame}



\begin{frame}

\begin{minipage}{.5\textwidth}


{\color{miugray}$\blacksquare ~$} \textbf{Aufgabe 4} \par \vspace*{2mm}

Durch Mischen von Orangensaft (Extraktgehalt: 12 °Brix, Dichte: 1,0465 kg/L, Gesamtsäure: 9
g/L, berechnet als Weinsäure WS), einer Saccharoselösung (c = 65\% w/w) und Wasser soll ein
Orangennektar mit einem Extraktgehalt von 8 °Brix (Dichte: 1,0299 kg/L) und einem
Gesamtsäuregehalt von 5 g/L (berechnet als WS) hergestellt werden. Ein Zusatz von
Säuerungsmitteln ist bei der Herstellung von Fruchtnektaren nicht üblich.

\begin{itemize}


\item Berechnen Sie, wieviel L Orangennektar aus 1 L Saft hergestellt werden können.



\end{itemize}
\end{minipage}%
\hspace*{4mm}
\begin{minipage}{.5\textwidth}

\antwort{
\begin{align*}
\uncover<2->{\frac{9~\text{g}}{\text{L}} \cdot 1 \text{L} &= \frac{5~\text{g}}{\text{L}} \cdot x  \\}
\uncover<3->{x &= \frac{9~\text{g}}{\text{L}} \cdot \frac{\text{L}}{5~\text{g}} \\}
\uncover<4->{x &= \frac{9}{5} = 1,8~\text{L}}
\end{align*}

}







\end{minipage}%


\end{frame}


\begin{frame}

\begin{minipage}{.5\textwidth}


{\color{miugray}$\blacksquare ~$} \textbf{Aufgabe 4} \par \vspace*{2mm}

Durch Mischen von Orangensaft (Extraktgehalt: 12 °Brix, Dichte: 1,0465 kg/L, Gesamtsäure: 9
g/L, berechnet als Weinsäure WS), einer Saccharoselösung (c = 65\% w/w) und Wasser soll ein
Orangennektar mit einem Extraktgehalt von 8 °Brix (Dichte: 1,0299 kg/L) und einem
Gesamtsäuregehalt von 5 g/L (berechnet als WS) hergestellt werden. Ein Zusatz von
Säuerungsmitteln ist bei der Herstellung von Fruchtnektaren nicht üblich.

\begin{itemize}



\item Berechnen Sie, welche Mengen an Saccharoselösung und Wasser entsprechend zu dem
Saft zugegeben werden müssen, um 1 kg des Getränkes herzustellen. Tipp: Berechnen Sie
dazu zunächst den Massenanteil an Fruchtsaft im Orangennektar.



\end{itemize}
\end{minipage}%
\hspace*{4mm}
\begin{minipage}{.5\textwidth}

\antwort{

\begin{align*}
\uncover<2->{m_1 &=}\uncover<3->{ 1,8 ~\text{L} \cdot 1,0299 ~\text{kg/L} = 1,854 ~\text{kg} \\}
\uncover<2->{m_2 &=}\uncover<3->{ 1 ~\text{L} \cdot 1,0465 ~\text{kg/L} = 1,0465 ~\text{kg} \\ }
\uncover<4->{&\rightarrow \frac{1,0465 ~\text{kg}}{1,854 ~\text{kg}} =  0,5644 ~\text{kg}}
\end{align*}
\uncover<5->{Um 1 Kg Nektar herstellen zu können benötigen wir 0,5644 kg Saft. \par }
\uncover<6->{Menge Zucker aus Saft berechnen:}
\uncover<7->{\begin{align*}
0,5644 ~\text{kg} \cdot 0,12 = 0,068 ~\text{kg}
\end{align*}}
\uncover<8->{Da wir 1 kg $\cdot $ 8° Brix = 0,08 kg Zucker in unserem Nektar benötigen, fehlen noch:
\begin{align*}
0,08~\text{kg} - 0,068~\text{kg} = 0,012 ~\text{kg}
\end{align*}}
\uncover<9->{Dafür benötigen wir folgende Menge 65\% Saccharose-Lösung:
\begin{align*}
\frac{0,012~\text{kg}}{0,65} = 0,0185 ~\text{kg}
\end{align*}}
\uncover<10->{Und der Rest ist Wasser:
\begin{align*}
(1-0,5644-0,0185)~\text{kg} = 0,4171 ~\text{kg}
\end{align*}}

}







\end{minipage}%


\end{frame}


















\begin{frame}

\begin{center}
\LARGE
Vielen Dank fürs Zuhören!
\end{center}

\end{frame}







\end{document}
